\scanpage{16}
% archive foliation: 15
\unsure{si on dispose d'une théorie analogue}
pour les complexes de chaînes de $k$-modules
ps. coh., ce qui \unsure{n'est guère le cas que}
si on suppose $k$ principal, (\unsure{Dedekind}
à la rigueur). Supposons donc
$k$ \deleted{\unsure{anneau}} \unsure{principal}
\deleted{\uline{principal}},
\inserted{(le cas $k$ \unsure{un corps} est déjà pratiquement le
\unsure{cas des $L_*$ spéciaux} \illegible{1} \unsure{précédent})}
et soit $\operatorname{Max}(k)$ l'ens.
des idéaux \deleted{premiers} \inserted{maximaux} de $k$,
$\unsure{\mathfrak P}=\mathbf N^{(\operatorname{Max}k)}-\{0\}$.
\margin{\unsure{la classe d'un projectif de t.f. est son rang}}

\uline{Th.} Tout $L_*$ \inserted{str.} $\partial$-ps. coh. est isomorphe
à un produit de modules de la
forme suivante
\[
\begin{array}{lll}
a)&\deleted{\unsure{M_i\otimes_k}}(\Lambda^i\Phi_*)^{\mu_i}
 &(i\geq1)\quad\deleted{M_i\ k\text{-module libre}}\ \inserted{\mu_i\in\mathbf N}\\
&&\deleted{\text{libre de t.f.}}\\
b)&\deleted{\unsure{N_j\otimes_k}}(\Lambda^j\Psi_*)^{\nu_j}
 &(j\geq0)\quad\deleted{N_j}\ (\nu_j\in\mathbf N)\\
c)&\deleted{\unsure{P_{i,\ell}\otimes_k}}\Phi(i,\ell)_*^{\beta_{i,\ell}}
 &(i\geq1)\quad\deleted{\unsure{P_{i,\ell}}}\ \beta_{i,\ell}\in\mathbf N,
\end{array}
\]
\deleted{où $\Phi(i,\ell)$}
$\ell$ un \deleted{idéal} diviseur de $k$ (cf plus bas
la définition de $\Phi(i,\ell)_*$), de façon
précise
\[
\begin{split}
L_*&\simeq\prod_{i\geq1}M_i\otimes_k\Lambda^i\Phi_*
\times\prod_{j\geq0}N_j\otimes_k\Lambda^j\Psi_*\\
&\qquad\times\prod_{\substack{i\geq1\\\ell\in\unsure{\mathfrak P}}}
\Phi_{i,\ell}^{\beta_{i,\ell}}
\end{split}
\]
où pour tout
$i\geq1$, l'ensemble des $\ell$ où les $\beta_{i\ell}$ sont non
nuls est fini. (NB de tels
produits sont bien str. $\partial$-ps coh.)
On peut supposer de plus que pour tout
$i$ donné, \unsure{l'ens.} des $\ell\in\unsure{\mathfrak P}$ \deleted{\unsure{les}} tels
que $\beta_{i\ell}\neq0$ est tot. ordonné, moyennant
